% ---------------------------------------------------------------------
%  §1. Resumen
% ---------------------------------------------------------------------
\section*{Resumen}

Este informe analiza cómo se asocia el uso de redes sociales con el sueño, la salud mental y el rendimiento
académico percibido de 1.743 estudiantes de 18 a 24 años de 110 países. Mediante un análisis descriptivo
—comparación de medias, medianas y proporciones por grupo, correlaciones y visualizaciones— se evalúan el
nivel de desarrollo del país, la plataforma más usada, el género y la situación sentimental. El impacto no es
igual para todos: los estudiantes de países en desarrollo, quienes declaran que las redes afectan su
rendimiento y los usuarios de WhatsApp, TikTok y Snapchat presentan perfiles menos favorables. El factor
común es la \textbf{cantidad de horas de uso}, que se correlaciona negativamente con el sueño ($-0{,}79$) y
con la salud mental ($-0{,}82$); el género y la situación sentimental casi no muestran diferencias.

\noindent\textbf{Palabras clave:} redes sociales, estudiantes, salud mental, sueño, rendimiento académico.
